
El capítulo anterior construye simultáneamente las cinco operaciones desde el mismo
estado enriquecido del continuo. Aquí se despliega cada componente completo antes de
su aplicación, con fibra, extensiones, relaciones, números, orientación, acarreo,
memoria, frontera, ruta, multisección, supervivencia, terminales y lectores. La
separación expositiva de las cinco tuplas no las convierte en ramas, factores ni
objetos independientes: son cinco características intrínsecas del único continuo
generado y conservan literalmente el mismo estado y el mismo ledger. El constructor
común no las reúne después; las produce simultáneamente como estructura interna de
un solo objeto.

\section{Dominio, mapas y compatibilidades del objeto terminal común}

Para exponer cada característica sin separarla de las otras, se usa el índice
$T\in\{\mathrm{sp},\mathrm{sol},\mathrm{gau},\mathrm{coh},\mathrm{det}\}$,
pero el dato matemático completo sigue siendo el mismo continuo enriquecido. La
vista interna no se postula de nuevo: es exactamente la lectura tipada construida en
\eqref{eq:pdf2-operacion-infinita-tipificada},
\begin{equation}\label{eq:pdf2-g-t-desde-continuo}
 \mathfrak G_T^{\rm int}:=
 \operatorname{View}_T\!\left(
 \mathcal O_\infty(\mathcal C_{\rm cont}^{\rm disc})\right).
\end{equation}
Su presentación completa en el canal \(T\) es
\begin{equation}\label{eq:pdf2-g-t-completo}
 \mathfrak G_T^{\rm int}=
 \bigl(
 \begin{gathered}
 \{\mathcal F_q^{(T)}\}_q,
 \{\operatorname{Ext}_\gamma^{(T)}\}_\gamma,
 \operatorname{Rel}_T,\operatorname{Num}_T,\operatorname{or}_T,\\
 \operatorname{Carr}_T,\operatorname{Mem}_T,\partial_T,\gamma_T,
 \mathfrak m_T,\operatorname{Surv}_T,X_{\chi_T},\tau_T,
 \operatorname{Term}_T,\operatorname{Lect}_T
 \end{gathered}
 \bigr).
\end{equation}
No se admite sustituir este objeto por $X_{\chi_T}$, por su cardinal o por el nombre
de un problema clásico. Sea $\mathcal D_T^{\rm cl}$ el espacio de datos clásicos
del problema correspondiente. Éstos entran sólo mediante
\begin{equation}\label{eq:pdf2-mapa-preparacion-aplicacion}
 \operatorname{App}_{T,d_T}:\mathfrak G_T^{\rm int}\longrightarrow
 \mathfrak G_T(d_T),\qquad d_T\in\mathcal D_T^{\rm cl},
\end{equation}
donde $\mathfrak G_T(d_T)$ es, por definición, el par formado por la vista interna
completa y las coordenadas de interfaz que se despliegan abajo. Así ninguna fibra,
relación, orientación, carry, memoria, frontera, ruta, multisección, supervivencia
o terminal se pierde al introducir el dato clásico. La proyección interna
$\operatorname{pr}_{X,T}:\mathfrak G_T^{\rm int}\twoheadrightarrow X_{\chi_T}$
es sobreyectiva. Si
$\mathcal A_{T,d_T}:\mathfrak G_T(d_T)\to\mathcal M_T$, se define
\[
 \widetilde{\mathcal A}_{T,d_T}:=
 \mathcal A_{T,d_T}\circ\operatorname{App}_{T,d_T},
 \qquad
 \mathfrak M_T:=\operatorname{Im}\mathcal A_{T,d_T}.
\]

\section{Despliegue interno de las realizaciones correlativas anterior a su identificación con problemas clásicos}

La fórmula genérica anterior no sustituye la escritura de las cinco características.
En este volumen cada interfaz queda materializada sin convertir la vista interna en
un objeto o dominio independiente. En los cinco displays siguientes, la primera
coordenada es siempre $\mathfrak G_T^{\rm int}$; por ello las listas de coordenadas
específicas amplían el registro común y nunca lo sustituyen.
Usaremos
\[
 d_{\rm sp}=(R,\mathcal D_R^{\rm cof}),\quad
 d_{\rm sol}=(u_0,f,T),\quad
 d_{\rm gau}=(\Lambda_m,g),\quad
 d_{\rm coh}=(X,p),\quad
 d_{\rm det}=E/\mathbb Q.
\]
Esta línea fija el instante causal exacto en que aparecen los problemas clásicos:
después de \(\mathfrak G_T^{\rm int}\), nunca dentro de su generación HMT.

\subsection{Acción espectral–polar del constructor común y conservación de la fibra crítica}

La característica espectral--polar conserva el registro
\begin{equation}\label{eq:pdf2-g-sp-explicito}
 \mathfrak G_{\rm sp}(d_{\rm sp})=
 \bigl(
 \begin{gathered}
  \mathfrak G_{\rm sp}^{\rm int};\\
  \widehat X_\infty,\Gamma_9,\widehat R=(R,K),
  \operatorname{Carr}_R,S_0,M,\mathcal O_9,\\
  \Xi_R,P_R,\mathcal D_R^{\rm cof},
  \mathcal F_R,\operatorname{Ext}_R,\operatorname{Surv}_R,
  \tau_R,\mathcal P_R
 \end{gathered}
 \bigr).
\end{equation}
Aquí $\mathcal D_R^{\rm cof}$ es el dominio de funciones de prueba; $\Xi_R$ publica
los canales primo, gamma y polar en un mismo Hilbert; $P_R=P_R^*=P_R^2$ es la
proyección natural; $S_0,M,\mathcal O_9$ conservan la memoria y el terminal; y
$\widehat R=(R,K)$ conserva tanto la ventana como su profundidad. El objeto de Weil
es el complemento, no un signo añadido:
\begin{equation}\label{eq:pdf2-g-sp-complemento}
 B_{W,R}:=\Xi_R^*(I-P_R)\Xi_R.
\end{equation}
Los canales regulados, la naturalidad cofinal y el terminal pertenecen a la misma
característica; ninguno se reconstruye desde los ceros de $\zeta$. La estructura
interna aporta la involución de hojas, la expectativa de extensiones, el defecto y
la telescopía terminal. El propietario exacto anterior construye además la
interfaz de Weil: fija la contracción de hojas en
\eqref{eq:amp-rh-owner-q}, liga el complemento espectral--polar en
\eqref{eq:amp-rh-owner-binding}, conserva el terminal mediante
\eqref{eq:amp-rh-owner-telescopia} y prueba su extinción en
\eqref{eq:amp-rh-owner-terminal}. En la notación del constructor común, esa
coligación ya construida se publica como
\begin{equation}\label{eq:pdf2-g-sp-coligacion-dato}
 \Sigma_R^{\rm sp}:
 \mathcal H_{+,R}\oplus\mathcal X_R^{\rm sp}
 \longrightarrow
 \mathcal H_{-,R}\oplus\mathcal L_R^{\rm sp}
 \oplus\mathcal X_R^{\rm sp},
 \qquad
 (\Sigma_R^{\rm sp})^*\Sigma_R^{\rm sp}=I,
\end{equation}
con estado terminal $x_{R,N}=q^NV_R^Nx_{R,0}$, $q=5/9$, y salida
$K_R^{\rm sp}J_{+,R}=J_{-,R}$. Su identidad de energía es
\begin{equation}\label{eq:pdf2-g-sp-identidad-finita-dato}
 \|J_{+,R}f\|^2
 =\|J_{-,R}f\|^2+
   \sum_{n=0}^{N-1}\|\ell_{R,n}f\|^2+
   q^{2N}\|x_{R,0}f\|^2.
\end{equation}
La factorización \eqref{eq:amp-rh-owner-l} y su publicación cofinal
\eqref{eq:amp-rh-owner-publicacion} prueban estas identidades y el
entrelazador primo--gamma--polar antes de este ensamblaje. El capítulo de
Riemann las despliega y reconoce mediante el
teorema~\ref{thm:amp-rh-cierre-cofinal-propietario}; no construye a
posteriori una interfaz ausente ni la recibe de los ceros de $\zeta$.

\subsection{Balance solenoidal del constructor común y conservación del flujo}

Para cada corte $M$ y profundidad $j$, la imagen es
\begin{equation}\label{eq:pdf2-g-sol-explicito}
 \mathfrak G_{{\rm sol},M,j}(d_{\rm sol})=
 \bigl(
 \begin{gathered}
  \mathfrak G_{\rm sol}^{\rm int};\\
  \Xi_{M,j}^{\rm sol},C_{M\leftarrow N},M_{M,j}^+,M_{M,j}^-,
  D_{M,j},H_{M,j},E_9,J_9,P_9,Q_{\rm micro},\\
  \Gamma_{M,j}^{\rm term},\mathcal E_{M,j,T}^{\rm term},
  L_{M,j,T}^{\rm term},\ell_{M,j},\operatorname{Carr}_{M,j},\\
  \operatorname{Mem}_{M,j},\operatorname{Surv}_{M,j},\tau_{M,j}
 \end{gathered}
 \bigr).
\end{equation}
El campo Galerkin publicado es una coordenada de $\Xi_{M,j}^{\rm sol}$, pero el
estado incluye además el defecto $C_{M\leftarrow N}$, las memorias orientadas, la
counidad y su sección, la dinámica terminal y la pérdida observada. El Gram
\begin{equation}\label{eq:pdf2-g-sol-gram-estructural}
 U_{M,j;J,T}^{\rm term}
 :=(\mathbf G_{M,j;J,T}^{\rm term})^*
    \mathbf G_{M,j;J,T}^{\rm term}
\end{equation}
conserva en $\mathbf G^{\rm term}$ la coordenada terminal y todas las pérdidas
propagadas. La telescopía no borra ninguna de ellas.
El propietario exacto que precede a este ensamblaje construye el mapa causal
hidrodinámico completo: el lift PC--Fock de
\eqref{eq:ns-import-pcf-lift}, la coligación isométrica de
\eqref{eq:ns-import-cascade-isometry}, su telescopía de Stein
\eqref{eq:ns-import-stein-telescope}, el ledger
\eqref{eq:ns-import-ledger}, el carry financiado una sola vez
\eqref{eq:ns-import-carry-funded} y el retorno exacto
\eqref{eq:ns-import-exact-return}. En esas mismas coordenadas, las ondas de
scattering se publican como
\begin{equation}\label{eq:pdf2-g-sol-ondas-dato}
 a_{M,j}=(2\sqrt\nu)^{-1}A_M^{-1/2}\Lambda^sf_{M,j},
 \qquad
 b_{M,j}=\sqrt\nu A_M^{1/2}\Lambda^su_{M,j}-a_{M,j},
\end{equation}
y la transición lineal actúa sobre todas las coordenadas mixtas
$u^{\odot r}\otimes f^{\odot k}$. La raíz común y su transporte no se
postulan en este bloque: ya están construidos por
\eqref{eq:nsclose-owner-metric-definition}--
\eqref{eq:nsclose-owner-root-uniformity}. Su publicación en la carta común es
\begin{equation}\label{eq:pdf2-g-sol-raiz-dato}
 J_T^{\rm sol}:\mathcal D_T^{\rm NS}\longrightarrow\mathcal H_T^{\rm sol},
 \qquad
 \mathbf G_{M,0;J,T}^{\rm term}\Lambda_{M,T}J_T^{\rm sol}
 =I_{M,J}J_T^{\rm sol},
 \qquad I_{M,J}^*I_{M,J}=I.
\end{equation}
La conservación conjunta de terminal, disipación, advección, carry,
microestado, shell y memoria está probada por el telescopio propietario
\eqref{eq:nsclose-owner-telescope} y su ledger
\eqref{eq:nsclose-proprietary-ledger}. El presupuesto uniforme, independiente
de $M,J$, es el teorema ya establecido por
\eqref{eq:nsclose-proprietary-uniform-budget} y
\eqref{eq:nsclose-uniform-sobolev-precompactness}; aquí se publica en la
notación del constructor común como
\begin{equation}\label{eq:pdf2-g-sol-cota-dato}
 \sup_{M,J}
 \|\mathbf G_{M,0;J,T}^{\rm term}\Lambda_{M,T}J_T^{\rm sol}d\|^2
 \le C_T^{\rm sol}\|d\|_{\mathcal D_T^{\rm NS}}^2,
\end{equation}
con $C_T^{\rm sol}$ independiente de $M,J$. El capítulo posterior despliega
la realización de Navier--Stokes y reconoce su codominio clásico; no aporta
una premisa ausente ni selecciona la cota desde la conclusión.

\subsection{Incidencia proyectiva de calibre del constructor común y descenso al cociente}

La imagen gauge no es la marginal de enlaces. Su objeto completo por nivel es
\begin{equation}\label{eq:pdf2-g-gau-explicito}
 \mathfrak G_{{\rm gau},m}(d_{\rm gau})=
 \bigl(
 \begin{gathered}
  \mathfrak G_{\rm gau}^{\rm int};\\
  \Lambda_m,\widehat X_m,\widehat\mu_m^{\rm ov},
  \{\widehat\mu_{m,g}^{\rm YM},\nu_{m,g},K_{m,g}^{\rm phys}\}_{g>0},\\
  \mathcal G_m^{\rm full},\mathbb P_m^{\rm phys},
  \widehat J_m,\widehat\Gamma_m,
  \{E_{m,c}\}_c,\widehat H_m^{\rm ov},D_m^{\rm YM},\\
  \{\Theta_{\nu,m}\}_{\nu=1}^4,
  \{\mathfrak A_{+,\nu,m},\operatorname{ev}_{t}^{(\nu)},
      \mathcal V_{m,\nu}^{\rm OS}\}_{\nu=1}^4,\\
  \widehat\Omega_m^{(8)},P_{\Omega_m},W_c,
  \operatorname{Carr}_m,\operatorname{Mem}_m,
  \operatorname{Surv}_m,\tau_m,\mathcal P_m^{F^2}
 \end{gathered}
 \bigr).
\end{equation}
El grupo $\mathcal G_m^{\rm full}$ actúa simultáneamente en enlaces, marcos y
coordenadas de incidencia $q_c\in\mathbb F_3^6$. La proyección física es su promedio
de Haar. El generador reducido es
\begin{equation}\label{eq:pdf2-g-gau-generador-estructural}
 D_m^{\rm YM}
 =\left.3\kappa_{\rm av}\sum_c(I-E_{m,c})
  \right|_{\operatorname{Ran}\mathbb P_m^{\rm phys}},
\end{equation}
y $W_c$ pertenece a esa subálgebra. Las cuatro reflexiones actúan sobre una misma
medida, un mismo kernel y un mismo terminal.
El propietario exacto anterior construye la acción gauge de incidencia
\eqref{eq:amp-ym-accion-gauge-incidencia}, la proyección física de Haar
\eqref{eq:amp-ym-proyeccion-fisica-haar} y la compresión entrelazada
\eqref{eq:amp-ym-compresion-entrelazada}. De esa construcción proceden el
kernel completo equivariable $\widehat K_{m,t}^{\rm ov}$, su descenso al
cociente físico $K_{m,t}^{\rm phys}$, los cuatro operadores de reflexión y
la acción $U_m(g)$ de $E(4)$ en el núcleo cilíndrico suave. Este bloque
ensambla, sin introducirlos como datos, el generador, el vacío, el testigo
de curvatura y las formas OS:
\begin{equation}\label{eq:pdf2-g-gau-entrelazamiento-dato}
 D_{m+1}^{\rm YM}\widehat J_m=\widehat J_mD_m^{\rm YM},
 \quad
 \mathbb P_{m+1}^{\rm phys}\widehat J_m
  =\widehat J_m\mathbb P_m^{\rm phys},
 \quad
 \widehat J_mW_{m,c}=W_{m+1,c}.
\end{equation}
La primera y la tercera identidades son la publicación, en la notación
común, de \eqref{eq:amp-ym-compresion-entrelazada} y
\eqref{eq:amp-ym-testigo-fisico-gap}; el semigrupo, el vacío simple y la
brecha exacta ya se obtienen en
\eqref{eq:amp-ym-owner-semigrupo-vacio}--
\eqref{eq:amp-ym-owner-brecha-exacta}. Aquí
$\nu_{m,g}:=(\pi_m^{\rm phys})_*\widehat\mu_{m,g}^{\rm YM}$ y
$K_{m,g}^{\rm phys}(t)=e^{-tD_{m,g}^{\rm YM}}$ en el cociente físico.
La completitud gauge incluye además la desintegración por cada familia de
hiperplanos euclídeos. Si $A_j\in\mathfrak A_{+,\nu,m}$ y
$t_1\le\cdots\le t_n$, la marginal de la misma medida física satisface
\begin{equation}\label{eq:pdf2-g-gau-transfer-euclidean-dato}
\begin{aligned}
 &\int\prod_{j=1}^n
   (\tau_{t_je_\nu}A_j)(x)\,d\widehat\mu_{m,g}^{\rm YM}(x)\\
 &\quad=
 \int\prod_{j=1}^nA_j(y_j)\,
  \nu_{m,g}(dy_1)
  \prod_{j=1}^{n-1}K_{m,g}^{\rm phys}
       (t_{j+1}-t_j;y_j,dy_{j+1})\\
 &\quad=
 \left\langle\mathbf1,
  M_{A_1}e^{-(t_2-t_1)D_{m,g}^{\rm YM}}M_{A_2}\cdots
  e^{-(t_n-t_{n-1})D_{m,g}^{\rm YM}}M_{A_n}\mathbf1
 \right\rangle.
\end{aligned}
\end{equation}
La reconstrucción anterior prueba la identidad en las cuatro direcciones
mediante \eqref{eq:amp-ym-cuatro-os} y construye un único Hamiltoniano
entrelazado en \eqref{eq:amp-ym-hamiltoniano-comun}. En consecuencia, la
traslación OS se publica aquí como
\begin{equation}\label{eq:pdf2-g-gau-os-generator-dato}
 \mathcal V_{m,\nu}^{\rm OS}T_{m,\nu}^{\rm OS}(s)
 (\mathcal V_{m,\nu}^{\rm OS})^{-1}
 =e^{-sD_{m,g}^{\rm YM}}.
\end{equation}
Así, el generador de expectativas y el Hamiltoniano de traslaciones no se
identifican por nombre: son dos realizaciones ya construidas de la misma
desintegración, cuya forma límite y cuya publicación de curvatura constan en
\eqref{eq:amp-ym-forma-limite} y \eqref{eq:amp-ym-f2}.

\subsection{Coherencia cohomológica del constructor común y emergencia de la clase global}

Para una variedad proyectiva lisa compleja $X$ y una codimensión $p$, la imagen
cohomológica conserva
\begin{equation}\label{eq:pdf2-g-coh-explicito}
 \mathfrak G_{\rm coh}(d_{\rm coh})=
 \bigl(
 \begin{gathered}
  \mathfrak G_{\rm coh}^{\rm int};\\
  \{\mathscr S_N(X,p),\rho_{N+1,N}\}_N,
  \{\mathscr H_N^p(X),q_{N+1,N}\}_N,
  \{e_N\}_N,
  \{\iota_{N+1,N}\}_N,\\
  \{\mathcal G_{N+1}^{\rm coh},\mathcal R_{N+1}^{\rm coh},
       D_{N+1},a_{N+1},b_{N+1}\}_N,\\
  \mathcal F_{X,p},\operatorname{Ext}_{X,p},X_\chi(X,p),
  \operatorname{Rel}_{X,p},\operatorname{Carr}_{X,p},
  \operatorname{Mem}_{X,p},\partial_{X,p},\gamma_{X,p},\\
  \Sigma_{X,p},\operatorname{Surv}_{X,p},\tau_{X,p},
  R_{\rm Hdg},\operatorname{ch}^{\rm loc}_p,\operatorname{cl}_{X,p}
 \end{gathered}
 \bigr).
\end{equation}
Los cuadrados
\begin{equation}\label{eq:pdf2-g-coh-cuadrados}
 e_N\rho_{N+1,N}=q_{N+1,N}e_{N+1}
\end{equation}
ligan observación, extensión y refinamiento. El propietario exacto anterior
construye, y no recibe como dato, la historia basada con sus correcciones
relativas, carry, memoria, frontera y terminal. En particular,
\eqref{eq:amp-hodge-mapas-relativos}--
\eqref{eq:amp-hodge-factorizacion-relativa} construyen los mapas del
incremento y \eqref{eq:amp-hodge-seccion-relativa} construye su inverso
derecho. En la notación del constructor común, ese mapa relativo se publica
como
\begin{equation}\label{eq:pdf2-g-coh-seccion-dato}
 s_{N+1}^{\rm rel}:\ker q_{N+1,N}\longrightarrow\ker\rho_{N+1,N},
 \qquad
 e_{N+1}^{\rm rel}s_{N+1}^{\rm rel}=I,
\end{equation}
natural bajo refinamiento por
\eqref{eq:amp-hodge-naturalidad-seccion-relativa} y compatible con carry y
terminal. La sección inicial tampoco se supone: se construye en
\eqref{eq:amp-hodge-seccion-base}. Son estos mapas ya demostrados, y no una
clase final elegida, los que el bloque ensambla y el capítulo posterior
reconoce en el codominio Hodge.

La completitud se apoya en acciones e identidades construidas por el mismo
propietario. La degeneración unitaria, la presentación relativa y el lector
se ensamblan aquí mediante
\begin{equation}\label{eq:pdf2-g-coh-datos-operativos}
\begin{gathered}
 \iota_{N+1,N}:\mathscr S_N(X,p)\longrightarrow\mathscr S_{N+1}(X,p),
 \qquad \rho_{N+1,N}\iota_{N+1,N}=I,\\
 \mathscr A_{N+1}^{\rm rel}
  =\mathbb Q^{(\mathcal G_{N+1}^{\rm coh})}/\operatorname{im}D_{N+1},
 \qquad e_{N+1}^{\rm rel}a_{N+1}=b_{N+1},
 \qquad b_{N+1}D_{N+1}=0.
\end{gathered}
\end{equation}
El conjunto $\mathcal G_{N+1}^{\rm coh}$ enumera exhaustivamente todos los
átomos nuevos y una etiqueta terminal por componente. La forma normal y el
inverso derecho se construyen por reducción basada en
\eqref{eq:amp-hodge-sigma-coordenada}; su compatibilidad con refinamientos se
prueba en \eqref{eq:amp-hodge-naturalidad-seccion-relativa}. Por ello, en el
orden por profundidad, soporte, hoja, ruta, carry, memoria y terminal, cada
átomo de observación tiene una única etiqueta máxima y
\begin{equation}\label{eq:pdf2-g-coh-tabla-dato}
 b_{N+1}(c_\mu)=\bar c_\mu+
   \sum_{\nu<\mu}B_{\nu\mu}\bar c_\nu,
 \qquad B_{\nu\mu}\in\mathbb Q.
\end{equation}
La triangularidad es una consecuencia de esa construcción basada, anterior a
cualquier clase $\alpha$; no es una hipótesis ni una tabla suministrada desde
el problema terminal.

La completitud y el lector del mismo objeto están tipados por
\begin{equation}\label{eq:pdf2-g-coh-limite-lector-dato}
 X_\chi(X,p)=\varprojlim_N\mathscr S_N(X,p),
 \qquad
 q_N\operatorname{ev}^{\rm Hdg}_{\infty,X,p}
  =e_N\operatorname{pr}_N,
 \qquad
 \operatorname{cl}_{X,p}\operatorname{ch}^{\rm loc}_pR_{\rm Hdg}
  =\operatorname{ev}^{\rm Hdg}_{\infty,X,p}.
\end{equation}
Las proyecciones $q_N$, incluida la coordenada terminal, separan puntos del
codominio de evaluación. La extensión operativa
\eqref{eq:amp-hodge-extension-operativa}--\eqref{eq:amp-hodge-ext}, el paso
al límite \eqref{eq:amp-hodge-xi-limite} y la publicación final
\eqref{eq:amp-hodge-publicacion-final} prueban ya que toda clase Hodge
arbitraria admite una historia compatible. El capítulo de aplicación
despliega y reconoce esa publicación; no suministra la existencia como una
premisa posterior.

\subsection{Línea determinantal del constructor común y evaluación aritmética terminal}

Para una curva elíptica $E/\mathbb Q$, la imagen determinantal conserva
\begin{equation}\label{eq:pdf2-g-det-explicito}
 \mathfrak G_{\rm det}(d_{\rm det})=
 \bigl(
 \begin{gathered}
  \mathfrak G_{\rm det}^{\rm int};\\
  \mathcal T_m^{\rm surv},G_{m,n},\Delta_{\rm AW},\varepsilon_{\rm AW},
  P_E^{\rm gen},1_E,P_E^K,P_E^{\rm PT},\\
  z_K,z_{\rm PT},X^{\rm orth},p_{\rm root},
  \mathcal C_E^K,\mathcal C_E^{\rm Iw},\mathcal C_E^{\rm PT},\\
  \ell_E,\mathcal C_E^{\rm loc,red},S_E(T),
  \operatorname{Carr}_E,\operatorname{Mem}_E,
  \operatorname{Surv}_E,\tau_E,\mathcal L_E
 \end{gathered}
 \bigr).
\end{equation}
La unidad $1_E$ se copublica en las ramas Kato y Poitou--Tate antes de la dualidad y
de la raíz. El producto fibrado conserva la sombra modular y la historia profunda.
El cono de localización, las páginas Bockstein y la línea determinante pertenecen a
una misma multisección basada; no son invariantes elegidos por separado después de
conocer el rango o el término principal.
El propietario exacto anterior construye la descomposición basada
$\Theta_E$, la homotopía del sumando finita--singular, el comparador completo
Kato--FERS, los mapas de páginas y gluing de Poitou--Tate y los mapas de
cociclos local--global. La cadena comienza en la línea determinante
\eqref{eq:amp-bsd-linea-determinante}, construye el mapa de cadenas
\eqref{eq:amp-bsd-phi-kato-fers}, el comparador de Poitou--Tate
\eqref{eq:amp-bsd-comparador-pt} y el triángulo local
\eqref{eq:amp-bsd-triangulo-localizacion}. Sus identidades operativas se
publican aquí como
\begin{equation}\label{eq:pdf2-g-det-identidades-dato}
\begin{gathered}
 \Phi_{K\leftrightarrow{\rm FERS}}\Psi_{K\leftrightarrow{\rm FERS}}=I,
 \qquad
 \Psi_{\rm PT}^{-1}\Psi_{\rm PT}=I,
 \qquad
 \Psi_{\rm PT}\Psi_{\rm PT}^{-1}=I,\\
 \ker\partial_{\rm loc}=\operatorname{im}\iota_{\Sha},
 \qquad
 \pi_\Phi\operatorname{Lift}_\Phi=I.
\end{gathered}
\end{equation}
Estas identidades son consecuencias del paquete exacto
\eqref{eq:amp-bsd-paquete-identidades}--
\eqref{eq:amp-bsd-paquete-composicion}. El comparador total se construye en
\eqref{eq:amp-bsd-comparador-total} y conserva la base común desde el lado
analítico hasta el aritmético; rango y término principal se publican después
en \eqref{eq:amp-bsd-igualdad-grados}--
\eqref{eq:amp-bsd-termino-principal}. El capítulo BSD expande y reconoce esa
composición sin seleccionar retrospectivamente ninguna de sus salidas.

\section{\(5\) lectores terminales no permutables del continuo genealógico}

El mapa de publicación $\mathcal P_T$ es el lector matemático que
traduce la estructura al codominio clásico. La cadena propia de este volumen es
\begin{equation}\label{eq:pdf2-cinco-cadenas}
 \mathfrak G_T^{\rm int}\xrightarrow{\operatorname{App}_{T,d_T}}
 \mathfrak G_T(d_T)\xrightarrow{\mathcal A_{T,d_T}}\mathfrak M_T
 \xrightarrow{\mathcal P_T}\mathcal C_T.
\end{equation}
Las correspondencias quedan fijadas una a una:
\[
 \mathcal C_{\rm sp}=\mathrm{RH},\quad
 \mathcal C_{\rm sol}=\mathrm{NS},\quad
 \mathcal C_{\rm gau}=\mathrm{YM},\quad
 \mathcal C_{\rm coh}=\mathrm{Hodge},\quad
 \mathcal C_{\rm det}=\mathrm{BSD}.
\]
En el ensamblado activo, los propietarios exactos que preceden a esta sección
construyen también las cinco publicaciones finales: RH en el
teorema~\ref{thm:amp-rh-cierre-cofinal-propietario}, Navier--Stokes en los
teoremas~\ref{thm:nsclose-proprietary-full} y
\ref{thm:nsclose-global-smooth}, Yang--Mills en el
teorema~\ref{thm:amp-ym-terminacion}, Hodge en el
teorema~\ref{thm:amp-hodge-completitud-coinductiva} y BSD en el
teorema~\ref{thm:amp-bsd-comparadores-construidos}. Esta sección registra su
ensamblaje correlativo en el único continuo genealógico; los capítulos
posteriores despliegan las cinco realizaciones y su reconocimiento clásico,
sin aportar premisas ausentes ni utilizar las conclusiones como entradas.

\begin{falsifier}[Ablación del dato de partida]
Si una demostración comienza en $X_{\chi_T}$, convierte una coordenada de
$\mathfrak G_T^{\rm int}$ en un selector opcional o sustituye una fibra por su cardinal, no ha
abreviado el mismo objeto: lo ha reemplazado. La conclusión no es transferible al
sustituto.
\end{falsifier}
